% ============================================================================
% CHAPTER 1: The Computers That Ask for Help
% Why even the smartest AI still needs humans — from Netflix pauses to 
% chatbots gone rogue
% ============================================================================
% Source: /markdown/ch1/Ch1_final.md
% ============================================================================

\begin{chapterquote}{The premise of this book}
Why even the smartest AI still needs humans---from Netflix pauses to chatbots gone rogue.
\end{chapterquote}

% ============================================================================
\section{The Pause That Makes You Think}
% ============================================================================

It's 2 AM. You've been binge-watching your favorite show for the past four hours, completely absorbed in the drama. The credits roll on yet another episode, and just as you're settling in for the next one, your screen goes dark except for a simple question:

\begin{center}
\textbf{``Are you still watching?''}
\end{center}

If you're like most people, this Netflix interruption probably annoys you. Maybe you've even googled ``how to turn off Netflix are you still watching'' in frustration. But here's the thing that might surprise you: this isn't Netflix being annoying. This is Netflix being \emph{smart}.

According to Netflix's own Help Center, the streaming giant asks this question ``so you don't lose your place or use internet data when you're not actually watching.'' The prompt appears after watching several episodes in a row without interaction, or after roughly 90 minutes of continuous viewing---the precise thresholds may vary by account type and region. Netflix's algorithm is constantly monitoring whether you've paused, skipped, or adjusted the volume---any sign that a human is still actively engaged.

That little pause isn't a bug---it's a feature. And it's one of the most common examples of something happening all around us, every single day, that most people never notice: computers asking humans for help.

% ============================================================================
\section{The Hidden Helpers Everywhere}
% ============================================================================

Once you start looking, you'll see these moments everywhere.

Your phone's autocorrect confidently changing ``I'm ducking tired'' instead of\ldots\ well, you know what you meant to type. The infamous ``duck'' substitution has become so ubiquitous that it's spawned countless memes, a Twitter account (@AvoidingDuck), and even made its way into academic papers about predictive text algorithms. Apple eventually addressed this in iOS~17, but for years, the system was so cautious about profanity that it created an entirely new vocabulary.

Siri responding with ``I'm not sure I understand'' when you ask her something in a noisy room. GPS navigation cheerfully announcing ``Recalculating'' when you've gone off its planned route. Your bank sending you a text asking ``Did you make this purchase?''

These aren't failures of technology. They're examples of something much more interesting: smart systems that are smart enough to know when they're confused.

Think about it. Netflix's algorithm could easily keep playing episodes forever. Your autocorrect could confidently change every word to whatever it thinks is best. Your voice assistant could always guess what you meant, even when it's not sure.

But they don't. And there's a very good reason why.

% ============================================================================
\section{When Smart Gets Smarter}
% ============================================================================

The smartest systems aren't the ones that never make mistakes---they're the ones that know when they might be making a mistake.

Netflix pauses your binge-watching session because its algorithms have learned something important: after several episodes or around 90 minutes of continuous watching without any interaction, there's a good chance you've fallen asleep, left the room, or gotten distracted. Rather than waste bandwidth streaming to an empty room (and potentially losing your place in the series), the system asks a simple question: ``Are you still there?''

It's a moment of uncertainty, handled gracefully.

Your voice assistant says ``I'm not sure I understand'' because it recognizes that the confidence level in its speech recognition has dropped below a safe threshold. Maybe there's background noise, maybe you mumbled, or maybe you asked about something it's never encountered before. Instead of guessing wrong and potentially doing something dangerous or embarrassing, it asks for clarification.

Your credit card company's fraud detection system works on the same principle. Banks like HSBC, Chase, and hundreds of others use machine learning algorithms that analyze transaction patterns in real-time. The system considers dozens of factors: the amount, the location, the merchant type, whether you've shopped there before, whether the purchase fits your normal spending patterns. When something looks unusual---say, a charge from a foreign country at an odd hour---the system doesn't just block the card or let it through. It sends you a text: ``Was this you?''

That two-way communication is the sweet spot: the algorithm handles the millions of obviously legitimate transactions automatically, but asks for human confirmation on the edge cases where it's genuinely uncertain.

% ============================================================================
\section{When the Asking Goes Wrong}
% ============================================================================

Of course, sometimes these systems get the asking part spectacularly wrong.

In January 2017, six-year-old Brooke Neitzel of Dallas, Texas, was chatting with her family's Amazon Echo when she said, ``Can you play dollhouse with me and get me a dollhouse?'' The device interpreted this as an order and charged the family \$170 for a KidKraft Sparkle Mansion dollhouse, along with four pounds of sugar cookies.

The story might have ended there as an amusing parenting lesson about voice-activated shopping. But when San Diego's CW6 News covered the story, anchor Jim Patton remarked on air: ``I love the little girl, saying `Alexa ordered me a dollhouse.'\,''

That sentence, broadcast to thousands of homes with their own Echo devices, set off a chain reaction. Viewers across San Diego reported that their own Echoes had woken up and attempted to order dollhouses. Amazon later clarified that their voice ordering system requires a confirmation step, but the incident revealed a fundamental problem: the system wasn't smart enough to distinguish between a child's casual conversation and an actual purchase request---or between a real command and a sentence spoken on television.

A few weeks later, during Super Bowl~LI on February~5, 2017, Google ran a commercial for its Home device. The ad repeatedly used the phrase ``OK Google'' as various actors demonstrated the product's features. Across America, viewers took to Twitter to report that their Google Home devices had activated and started trying to respond to the commercial's commands.

``My Google Home was confused during the Google commercial\ldots\ I felt bad,'' one viewer tweeted. ``Google Home ad on Super Bowl just freaked out my Google Home. Started trying to deal with queries as actors said,'' reported another.

The devices mostly responded with ``Sorry, I don't understand''---which, ironically, was actually the system behaving correctly. It recognized the commands but couldn't make sense of the context, so it asked for clarification. But the incident highlighted how difficult it is to build systems that know the difference between ``someone is talking to me'' and ``someone is talking about me on TV.''

% ============================================================================
\section{The Thermostat That Cried Wolf}
\label{sec:nest-thermostat}
% ============================================================================

Not all smart systems ask for help with the grace of Netflix. Some ask so badly that users wish they'd never tried to be smart at all.

Consider the Nest Learning Thermostat, Google's flagship smart home device. Introduced in 2011 with great fanfare, the Nest promised to learn your temperature preferences and create automatic schedules, saving energy while keeping you comfortable. One of its signature features is ``Home/Away Assist''---the thermostat uses sensors and your phone's location to detect whether anyone is home, automatically adjusting the temperature to save energy when you're out.

In theory, this is exactly the kind of \term{human-in-the-loop} design we've been celebrating. Instead of running heating or cooling constantly, the system asks itself: ``Is anyone here?'' and acts accordingly.

In practice, it's been a disaster for many users.

Google's own community forums are filled with complaints from frustrated homeowners. ``The Nest keeps thinking I'm away when I'm actually home,'' writes one user. ``I'll be sitting in my living room and suddenly the house starts cooling down because the thermostat decided nobody was here.'' Another reports: ``My Nest switches to Away mode multiple times a day while I'm working from home. I've basically disabled all the `smart' features and just use it as a dumb thermostat now.''

The problem is a perfect illustration of what happens when a system asks for help at the wrong times, with the wrong confidence thresholds, and with poor ability to learn from its mistakes.

The Nest relies on multiple signals: motion sensors in the device itself, your phone's GPS location, and activity patterns it has learned over time. But each of these can fail in predictable ways. Motion sensors can't see through walls, so if you're in a bedroom or home office, the thermostat in the hallway might not detect any activity. Phone location tracking is imprecise and battery-intensive. And learned patterns break down when your routine changes---say, when you start working from home, or when guests visit.

The key failure is that the Nest acts on its conclusions even when they're uncertain. When it decides you're ``away,'' it doesn't ask for confir\-mation---it just starts changing the temperature. There's no notification saying ``I haven't detected any motion for a while. Should I switch to energy-saving mode?'' No opportunity for you to say ``No, I'm actually here, I'm just being quiet.''

Compare this to Netflix. When Netflix thinks you might have stopped watching, it doesn't delete your place in the series and start playing something else. It simply pauses and asks. The cost of being wrong is minimal---a moment of annoyance for an engaged viewer. The Nest, by contrast, makes you too cold or too hot, wastes energy re-adjusting when it realizes its mistake, and generally teaches users that ``smart'' features are more trouble than they're worth.

Many Nest owners have shared a rueful common endpoint: they disabled Home/Away Assist entirely and went back to manual scheduling---effectively turning their \$250 smart thermostat into a \$30 programmable one. One tech reviewer titled their piece: ``I Love My Nest, But I Had to Turn Off Its Brain.''

The lesson here isn't that smart thermostats are a bad idea. It's that the \emph{way} a system asks for help---or the way it acts on uncertainty---matters as much as whether it asks at all. Netflix asks gracefully, with low stakes for being wrong. The Nest acts silently and confidently, even when its confidence is misplaced.

% ============================================================================
\section{The Systems That Didn't Ask}
% ============================================================================

If asking for help badly creates frustration, \emph{not} asking for help when you should can create tragedy.

GPS navigation systems have repeatedly directed drivers into dangerous situations because the algorithms were confident but wrong. In March 2021, a man in Charlton, Massachusetts, drove his car directly into Buffumville Lake after following his GPS navigation, which had apparently directed him down a boat launch ramp. The car was completely submerged in about 8--10 feet of water; the driver managed to escape and was transported to a hospital.

``The driver told rescuers he was following his GPS when his car suddenly plunged into the lake,'' the Charlton Fire Department reported.

Similar incidents have occurred repeatedly around the world. In May 2023, tourists in Hawaii followed their GPS right into Honokohau Harbor while looking for a manta ray tour. The viral video showed bystanders watching in disbelief as the car floated off the boat ramp. ``I didn't hear any brake sound or anything,'' a witness told local news. ``She didn't try to stop at all.''

In 2016, a woman in Ontario, Canada drove her car into Georgian Bay---an inlet of Lake Huron---after following GPS directions onto a boat launch in Tobermory \autocite{cbc2016gps}. She was rescued by locals who heard her honking for help.

These aren't failures of the GPS hardware---they're failures of system design. The navigation algorithms were confident in their directions but lacked the contextual understanding to recognize that a boat ramp is not a road, or that continuing forward would mean driving into water. A human, seeing the situation, would immediately recognize something was wrong. The system had no mechanism to express uncertainty, no way to say ``this doesn't look right.''

\begin{keyconcept}[Key Insight]
The most dangerous technology isn't the kind that asks for help too often---it's the kind that never asks at all.
\end{keyconcept}

% ============================================================================
\section{The High-Stakes Version}
% ============================================================================

The same principle that makes Netflix pause your binge-watching session applies in much more serious situations---but with much higher stakes for getting it wrong.

When Tesla's Full Self-Driving system encounters a scenario it's not confident about, it alerts the human driver to take control. According to Tesla's owner's manual, ``If Full Self-Driving (Supervised) disengages, an alert will appear on the touchscreen to notify you and a chime will sound. If this happens, take control of steering immediately.''

The system uses multiple layers of confidence monitoring: cameras watching the road, cameras watching the driver, sensors measuring hand pressure on the steering wheel. When any of these signals suggest uncertainty ---an unusual road situation, a distracted driver, sensor occlusion--- the system escalates, ultimately returning control to the human.

This isn't the system admitting defeat. It's the system being responsible about its limitations.

Medical AI systems work on similar principles. In radiology, AI algorithms are increasingly used to analyze X-rays, CT scans, and MRIs. But rather than simply giving a diagnosis, these systems are designed to express their confidence level. A 2024 article in \emph{Nature Machine Intelligence} described systems that provide \term{uncertainty quantification} in medical imaging---essentially, AI that tells the doctor ``I'm 95\% confident this is a fracture'' versus ``I'm only 60\% confident, so you should take a closer look'' \autocite{begoli2019uncertainty}.

Research has shown that AI systems that flag uncertain cases for human review serve as a ``safety net that may reduce misdiagnosis and improve patient outcomes.'' The key insight is that the AI doesn't need to be perfect---it just needs to know when it's not confident, and to route those cases to human experts.

% ============================================================================
\section{When Chatbots Pretend to Know}
\label{sec:air-canada}
% ============================================================================

In February 2024, a Canadian tribunal handed down a ruling that would send shockwaves through the tech industry \autocite{moffatt2024}. Jake Moffatt, grieving the loss of his grandmother, had visited Air Canada's website to book last-minute flights. He asked the airline's AI chatbot about bereavement fares---the discounted rates many airlines offer for emergency family travel.

The chatbot told Moffatt he could book a regular ticket now and apply for a partial refund to the bereavement rate within 90 days.

This was completely wrong. Air Canada's actual policy required passengers to apply for bereavement rates \emph{before} booking, not after. But the chatbot presented its incorrect information with the same breezy confidence it would use to list departure times or baggage fees. Moffatt, trusting the official airline website, booked his tickets accordingly.

When he later applied for the bereavement discount, Air Canada denied his request and pointed him to the official policy page---a different part of the same website that contradicted what the chatbot had told him.

Moffatt sued. Air Canada's defense was remarkable: the airline argued that it could not be held liable for information provided by its chatbot, essentially suggesting the chatbot was ``a separate legal entity that is responsible for its own actions.''

The tribunal didn't buy it. ``This is a remarkable submission,'' wrote Tribunal Member Christopher Rivers. ``While a chatbot has an interactive component, it is still just a part of Air Canada's website. It should be obvious to Air Canada that it is responsible for all the information on its website. It makes no difference whether the information comes from a static page or a chatbot.''

Air Canada was ordered to pay Moffatt \$812.02, including the fare difference, interest, and tribunal fees.

The case highlights a critical problem with AI systems: they can be ``confidently wrong.'' Studies across multiple domains have found that AI systems produce incorrect information---sometimes called \term{hallucinations}---at rates that vary widely depending on the task and system, but that remain significant even in purpose-built deployments. A 2024 study found that AI legal research tools hallucinated in over 58\% of legal queries \autocite{magesh2024hallucination}.

\begin{cautionbox}[The Air Canada Lesson]
The Air Canada chatbot's failure wasn't that it didn't know the bereavement policy. Its failure was that it didn't know that it didn't know---and had no mechanism to say ``I'm not certain about this, let me connect you with a human.''
\end{cautionbox}

% ============================================================================
\section{The Goldilocks Problem}
\label{sec:goldilocks}
% ============================================================================

Building AI that asks for help at the right moments is harder than it might seem. It's a classic ``Goldilocks problem''---not too much, not too little, but just right.

Ask for help too often, and you create \term{alert fatigue}. Think about websites that constantly ask if you want to enable notifications, cookie consent banners that appear on every single page, or security systems that cry wolf so frequently that people start ignoring them. Research in healthcare has shown that when clinical decision support systems generate too many alerts, physicians begin to override them automatically---sometimes missing the rare alerts that actually matter \autocite{lee2004trust}.

Ask for help too rarely, and you get dangerous overconfidence. This is what happened with some early autonomous vehicle systems that were so confident in their abilities that they didn't alert human drivers soon enough when they encountered situations beyond their capabilities. It's what happened with the Air Canada chatbot that didn't know to say ``I'm not sure.''

And then there's the Nest thermostat problem: asking in the wrong \emph{way}. The Nest technically ``asks'' whether you're home through its sensors and location tracking. But it doesn't communicate that question to you, and it acts on ambiguous answers as if they were certain. It's asking for help, but not in a way that actually involves you in the decision.

The goal is to find the sweet spot: systems confident enough to handle the vast majority of situations independently, but wise enough to recognize their limits and ask for human guidance when it matters most---and humble enough to ask in ways that actually help.

% ============================================================================
\section{What Makes a System ``Smart Enough to Ask for Help''?}
\label{sec:five-principles}
% ============================================================================

Looking across all these examples---from Netflix to self-driving cars to medical imaging---several key principles emerge:

\paragraph{Self-awareness.} The system needs to have some measure of its own confidence or uncertainty. It can't just make decisions; it needs to know how confident it is about those decisions. This is technically called ``uncertainty quantification,'' and it's one of the most active areas of AI research \autocite{kendall2017uncertainties}.

\paragraph{Context understanding.} The system needs to understand not just what it's being asked to do, but why it matters. The stakes involved in autocorrecting a text message are different from the stakes involved in medical diagnosis. Netflix can afford to occasionally ask ``Are you still watching?'' when you're actively engaged---the cost is a moment of annoyance. A self-driving car can't afford to be wrong about whether there's a pedestrian in the road.

\paragraph{Human psychology.} The system needs to ask for help in ways that humans can understand and respond to effectively. A cryptic error code doesn't help; a clear question does. This is why your bank doesn't text you ``Transaction flagged: anomaly score 0.73, confidence interval 0.45--0.89.'' It asks: ``Did you spend \$500 at Bangkok Electronics? Reply YES or NO.''

\paragraph{Graceful degradation.} When the system isn't sure, it needs to fail in a way that makes things better, not worse. Netflix pausing is better than streaming forever to an empty room. A GPS saying ``Recalculating'' is better than confidently directing you into a lake. An autocorrect that leaves a misspelling alone is often better than ``correcting'' it wrongly.

\paragraph{Appropriate escalation.} The Nest's failure wasn't that it tried to detect whether anyone was home. It was that when it was uncertain, it escalated in the wrong direction---acting silently on its conclusion instead of checking with the human. Good systems match their level of autonomous action to their level of confidence.

These five principles correspond to what this book calls the \textbf{Five Dimensions of \hitl System Design}:

\begin{keyconcept}[The Five Dimensions of HITL System Design]
\begin{description}
  \item[1.\ Uncertainty Detection] The system must recognize when it might be wrong---calibrated confidence, not false certainty.
  \item[2.\ Intervention Design] The ask must be designed for the human: well-framed, timely, and formatted for accurate response.
  \item[3.\ Timing] Human involvement must be triggered at the right moment---not too early, not too late, not so often that it becomes noise.
  \item[4.\ Stakes Calibration] The threshold for asking should reflect the cost of each type of error, not some generic accuracy target.
  \item[5.\ Feedback Integration] The human's response must flow back into the system---closing the loop so the system improves over time.
\end{description}
Each subsequent chapter deepens one or more of these dimensions. Chapter summaries throughout the book include a Five Dimensions Check to make the connections explicit.
\end{keyconcept}

% ============================================================================

\begin{trythis}
Over the next 24 hours, pay attention to how many times a piece of technology ``asks'' you something. Your phone, your streaming services, your apps, your smart home devices. Keep a quick count. Then ask yourself: Which of these felt helpful? Which felt annoying? What made the difference? You might be surprised how often computers are asking for your help---and how rarely we notice the ones that do it well.
\end{trythis}

% ============================================================================
\section{A Brief Technical Interlude}
\label{sec:technical-interlude}
% ============================================================================

You might be wondering: how does a computer program \emph{actually} know when it's uncertain? After all, computers are deterministic machines---they do exactly what they're programmed to do. How can they be ``unsure''?

The answer involves probability. Modern AI systems don't just produce answers; they produce \emph{probabilities} of different answers.

When your phone's autocorrect suggests ``duck'' for what you typed, it's not because it's certain you meant to write ``duck.'' It's because the algorithm calculated that ``duck'' had, say, a 60\% probability of being correct, while the\ldots\ alternative word\ldots\ had only a 40\% probability (because Apple's training data deliberately down-weighted profanity).

The system makes a decision based on which option is more probable. But it \emph{also} knows how close the decision was. When the top choice has 95\% probability, the system can be confident. When the top choice has only 51\% probability, the system should be uncertain---and ideally, should express that uncertainty by asking the user to confirm.

This same principle applies across all the examples we've discussed. Netflix's ``Are you still watching?'' appears when the probability that you're still actively engaged drops below some threshold. Fraud detection systems alert you when a transaction's probability of being legitimate falls into an uncertain range. Medical AI systems flag cases where their confidence level falls below the threshold for automatic diagnosis.

The technical challenge is \term{calibrating} these systems so that their expressed confidence actually matches reality. A system that says it's ``90\% confident'' should be right about 90\% of the time---not 70\% or 99\%. Getting this calibration right is what separates useful AI from dangerous AI \autocite{guo2017calibration}.

\begin{technote}[For the Curious Reader]
For readers interested in the technical details, the appendix to this chapter covers uncertainty quantification methods including Bayesian neural networks \autocite{gal2016dropout}, Monte Carlo dropout, deep ensembles \autocite{lakshminarayanan2017simple}, and calibration metrics.
\end{technote}

% ============================================================================
\section{The Future Is Collaborative, Not Automated}
% ============================================================================

This shift toward systems that know when to ask for help represents something bigger than just better technology---it represents a new philosophy about the relationship between humans and machines.

For a long time, the implicit goal of AI development seemed to be elimination: eliminate human error, eliminate human involvement, eliminate the need for human judgment entirely. The ultimate success would be full automation.

But the examples we've looked at suggest a different path: not replacement, but collaboration. The best systems aren't the ones that never need human help---they're the ones that know exactly when human help will make things better.

Netflix doesn't eliminate the need for human judgment about what to watch; it makes that judgment easier and more informed. Medical AI doesn't replace doctors; it augments their expertise with pattern recognition capabilities that exceed human perception, while flagging the cases that need human attention \autocite{kompa2021second}. Fraud detection systems don't replace your judgment about your own spending; they add a layer of protection that catches the things you might miss.

The Air Canada case shows what happens when this philosophy is ignored---when a system is designed to appear confident even when it shouldn't be, with no mechanism for admitting uncertainty or escalating to a human. The GPS-into-lake incidents show what happens when systems lack even basic sanity checks against their recommendations. And the Nest thermostat shows what happens when systems ``ask'' in ways that don't actually involve humans in the decision.

% ============================================================================
\section{Why You Should Care}
% ============================================================================

You might be thinking: ``This is interesting, but why does it matter to me? I'm not building AI systems.''

But here's the thing: you're already part of these systems. Every time you respond to Netflix's ``Are you still watching?'' you're participating in a human-AI collaboration. Every time you accept or reject an autocorrect suggestion, you're teaching the system about your preferences. Every time you confirm or deny a fraud alert, you're helping to train better security systems.

Understanding how these systems work---and why they ask for help---makes you a better collaborator. It helps you recognize when a system is being appropriately cautious versus when it's being unnecessarily annoying. It helps you provide better feedback when systems ask for your input. And it helps you recognize when a system might be confidently wrong, so you know when to double-check.

More importantly, as AI systems become more prevalent in everything from healthcare to education to transportation, we all have a stake in ensuring they're designed with appropriate humility. Systems that know their limits and ask for help when they need it are safer, more trustworthy, and ultimately more useful than systems that pretend to be infallible.

The tribunal ruling in the Air Canada case explicitly rejected the idea that chatbots should be treated as separate entities responsible for their own mistakes. Companies are responsible for the information their AI systems provide. As users, we have a right to expect that when systems are uncertain, they'll tell us---rather than presenting guesses as facts.

% ============================================================================
\section{What's Next}
% ============================================================================

This principle---smart systems knowing when to ask for help---is just the beginning. In the next chapter, we'll explore how to recognize those moments when smart systems get confused: What does it look like when an AI is uncertain? How do the best systems communicate their limitations? And why is that uncertainty not a failure, but an opportunity?

The future won't be about machines that never need help. It will be about machines that are smart enough to know exactly when they do need help, and wise enough to ask for it in ways that make both the human and the machine more effective.

That Netflix pause isn't interrupting your binge-watching session. It's giving you a glimpse of what thoughtful AI design looks like: systems that are powerful enough to work independently most of the time, but humble enough to check in when it matters.

The computers that ask for help aren't broken. They're actually the smartest ones of all.

\bigskip
\begin{center}
\rule{0.5\textwidth}{0.4pt}
\end{center}
\bigskip

\noindent\textit{In the next chapter, we'll explore how to recognize those moments when smart systems get confused---and why that confusion isn't failure, but opportunity.}

% ============================================================================
\section*{Chapter Summary}
\addcontentsline{toc}{section}{Chapter Summary}
% ============================================================================

\begin{keyconcept}[Key Concepts]
\begin{itemize}
    \item Smart systems that ask for help aren't malfunctioning---they're being intelligent
    \item Examples range from everyday (Netflix, autocorrect, fraud alerts) to high-stakes (medical AI, autonomous vehicles)
    \item The best AI systems know their confidence levels and recognize their limits
    \item Systems can fail in three ways: asking too often (alert fatigue), never asking (dangerous overconfidence), or asking badly (the Nest thermostat problem)
    \item ``Asking for help'' moments require sophisticated engineering, including uncertainty quantification
\end{itemize}
\end{keyconcept}

\subsection*{Key Examples}

\begin{description}
    \item[Netflix ``Are you still watching?''] System monitors engagement and asks after 3 episodes or 90 minutes; low-stakes, graceful design
    \item[Nest thermostat Home/Away Assist] System tries to detect occupancy but acts on uncertainty without confirming; frustrated users disable ``smart'' features
    \item[Amazon Echo dollhouse incident (2017)] System couldn't distinguish child's conversation from purchase request
    \item[Google Home Super Bowl activation (2017)] Devices activated by TV commercial, couldn't distinguish context
    \item[GPS into water incidents] Systems that never express uncertainty about directions
    \item[Air Canada chatbot case (2024)] Chatbot was ``confidently wrong'' about bereavement policy; tribunal ruled airline responsible
\end{description}

\subsection*{Key Principles}

\begin{itemize}
    \item The future is about human-AI collaboration, not full automation
    \item Well-designed systems are calibrated so that expressed confidence matches actual accuracy
    \item How a system asks for help matters as much as whether it asks
    \item Understanding these systems makes us better partners in human-AI teams
\end{itemize}
